\section{\label{III_3_B}Qualité des données : une tâche à part entière}

Ainsi, le contrôle de la cohérence des données au sein du modèle de données s'impose comme une tâche à part entière, à laquelle s'ajoute le contrôle de la qualité des données elles-mêmes. Si un pipeline de validation permet d'encadrer la création des données descriptives, les erreurs peuvent persister et dépendent du temps dédié à la validation. Dès lors, l'évaluation de la qualité des données produites s'ajoute à la vérification de leur cohérence structurelle au sein du modèle FRBR. 

À ce stade du processus, la qualité des données ne désigne pas de métrique technique mais plutôt une valeur métier, en lien avec les normes descriptives en vigueur. Deux niveaux d'exigence de qualité peuvent être établis. Premièrement, les données produites doivent être évaluées en lien avec les normes de description employées au sein de l'espace de travail. Par exemple, S-Movies est fondé sur les règles de description des archives (\gls{rda}) conformément aux reccommandations de la \gls{fiaf}. Deuxièmement, les données pourront être vérifiées en accord avec les règles internes et propres à chaque institution patrimoniale. Par exemple, les thésaurus, les autorités, les dates. Comme nous l'avions mentionné, les mots-clés produits sur le corpus étudié en seconde partie sont contrôlés par un thésaurus "keyword". Il s'agira alors de vérifier si les mots-clés versés dans l'application existent bien sous une forme contrôlée au sein d'un thésaurus dédié. 

Toutefois, ces vérifications laissées à l'utilisateur métier peuvent prendre un temps considérable et ne permettent pas de couvrir l'entiereté de la masse de données produites automatiquement. Dans ce cadre, l'usage de grands modèles de langage (\gls{llm}) peut être envisagé afin de contrôler la qualité des données au niveau institutionnel. En effet, Skinsoft développe déjà l'intégration d'un \gls{llm} pour les fonctionnalités de recherche en langage naturel et de \gls{rag}, entraîné sur les règles d'indexation métier générales et spécifique aux institutions. Le même \gls{llm} pourra être utilisé afin de pointer les incohérences au sein des notices contenant des données indexées automatiquement et manuellement. Un tel outil permettrait également d'homogénéiser l'ensemble des données indexées. Ainsi, l'utilisation d'un outil basé sur l'IA supplémentaire permettrait d'améliorer la production de données homogènes en regard des normes en vigueur, contrairement aux données indexées manuellement, souvent hétérogènes.  

Utiliser des LLM pour corriger les données et exécuter ces tâches de validation d'une manière hybride entre l'humaine et la machine, l'IA ne dégrade pas forcément la qualité des métadonnées mais peut l'améliorer en produisant des métadonnées homogènes au regard des normes. Selon Nikki Wise, certaines études montrent même que les utilisateurs préfèrent "les métadonnées générées par des machines ou par des systèmes hybrides à celles créées par des humains, en invoquant leur contexte plus riche, leur spécificité et leur exhaustivité [Traduction personnelle]" \footnote{\textit{"human users preferred machine-and hybrid-generated metadata over human-generated content, citing its increasedcontext, specificity, and comprehensiveness"} \cite{wise2026}}." Les capacités d'interprétation et d'analyse des \gls{llm} s'avèrent particulièrement efficaces pour assister l'utilisateur, et ce, même avec un entraînement limité, offrant ainsi "une perspective prometteuse pour les flux de travail hybrides associant humains et IA dans la création et la correction des métadonnées. [Traduction personnelle]" \footnote{\textit{"suggest promise for hybrid human-AI workflows in metadata creation and remediation"} \cite{wise2026}}.
Ce contrôle de la qualité des données ne peut toutefois pas s'effectuer sans leur traçabilité, prévue par les métadonnées de processus décrites précédemment, permettant d'avoir accès à l'historique des modifications et au contexte de création de ces données.

%Contrôle qualité et apprentissage continu : erreurs identifiées comme données d'annotation 
%Enfin, cette traçabilité alimente une boucle d'apprentissage continu. Les erreurs détectées par l'archiviste lors du contrôle qualité ne sont pas seulement corrigées dans la notice : elles sont requalifiées en données d'annotation. Réinjectées dans le pipeline, ces corrections enrichissent le jeu d'entraînement des modèles, perfectionnant ainsi les performances des futures propositions automatiques. 